\section*{Your turn}

Every chapter from here ends with something to do, and from
chapter~\ref{ch:hotchocolate-quickly} onward that means checking out a tag from
the companion repository and running it. This one needs no code, because the
question this chapter raises is about your organisation rather than your
toolchain.

Pick a graph. If your team runs a GraphQL API, use that one, however small. If
you do not have one to hand, use GitHub's, which publishes its whole schema as a
single downloadable file~\autocite{github2026schema} alongside the
breaking-change policy that keeps it
manageable~\autocite{github2026breaking}. It is large, genuinely load bearing,
and deliberately not federated, which makes it a useful control case for
everything argued above. Answer the questions you can from the schema and the
policy; where you cannot, note what a maintainer would have to tell you.

\begin{enumerate}
  \item \textbf{Map the owners.} List the fields on \code{Query}. Next to each
  one, write the name of the team that would review a change to it. Count the
  fields where you cannot name a team, or where you can name three.

  \item \textbf{Measure the blast radius.} Take the type that appears in the
  most operations. How many separate features read it? If you added a required
  argument to one of its fields tomorrow, who finds out, and how?

  \item \textbf{Time a change.} Find the last field added to the schema. From
  the first commit to serving production traffic, how long did it take, and how
  many people had to agree? Separate the time spent building from the time
  spent waiting.

  \item \textbf{Test the coupling.} Can one domain in that schema ship a change
  to production without the others shipping too? If the answer is no, is that
  because of the schema, the deployment, or the database?

  \item \textbf{Count the queues.} Over the last quarter, how often did two
  teams need to change the same type in the same week?
\end{enumerate}

If your answers are dull, that is the most useful outcome available to you
here. Fields with one obvious owner, a change that ships in a day, one
deployable that everybody is happy to release: you have a healthy single graph
and no federation problem. Part~I of this book will still be worth your time,
and you can treat the rest as a map of where the road goes.

Keep the notes either way. Chapter~\ref{ch:do-you-actually-need-federation}
turns this audit into an actual decision, and by then you will have seen what
the alternative costs.
